\section{Каузальность, локальность, унитарность}

В этой секции $\caStepSpace$ и~$\htick$ --- шаг и такт из определения элементарных шага и такта (\cref{def:step-tick}),
$\caSignalSpeed=\caStepSpace/\htick$ --- тактовая скорость сигнала на~$\layerMicro$,
$N(x)$ --- узлы, достижимые из~$x$ за один такт (окрестность одного такта, \cref{def:neighborhood}).

\begin{axiom}[A1 · Каузальность]
\label{ax:A1}
За один такт~$\htick$ сигнал проходит не дальше одного шага~$\caStepSpace$.
$N(x)$ --- узлы, до которых он доходит за этот такт.
\end{axiom}

\begin{axiom}[A2 · Локальность]
\label{ax:A2}
Закон $g$ локален: новое значение $z(x,t+1)$ зависит только от поля в $\varepsilon$-окрестности $x$:
\[
  z(x,t+1)=g\bigl(\{z(y,t)\}_{y\in N(x)}\bigr).
\]
\end{axiom}

\begin{figure}[htbp]
  \centering
  \begin{tikzpicture}[carrier panel]
  \CarrierHexRing{0.38}{0.65}
  \node[carrier site] at (0,0) {};
  \node[anchor=west, inner sep=2pt] at (0.06,0) {$x$};
  \node[anchor=south, inner sep=2pt] at (0,0.95) {$N(x)$};
  \draw[carrier arrow, carrier muted] (0.38,0) -- (0.65,0);
  \node[carrier muted, font=\scriptsize, anchor=south] at (0.52,-0.1) {$\caStepSpace$};
\end{tikzpicture}

  \caption{$A_1$: световая окрестность $N(x)$, $\caSignalSpeed\htick=\caStepSpace$.}
  \label{fig:axiom-a1-neighborhood}
\end{figure}

\begin{figure}[htbp]
  \centering
  \begin{tikzpicture}[carrier panel, x=0.75cm, y=0.75cm]
  \node[draw, carrier object, carrier fill, circle, minimum size=0.65cm, font=\scriptsize] (z1) at (0,0) {$z_1$};
  \node[draw, carrier object, carrier fill, circle, minimum size=0.65cm, font=\scriptsize] (z2) at (0,2) {$z_2$};
  \node[draw, carrier object, carrier fill, circle, minimum size=0.65cm, font=\scriptsize] (z3) at (1,1) {$z_3$};
  \node[draw, carrier object, carrier fill, circle, minimum size=0.65cm, font=\scriptsize] (z4) at (2,0) {$z_4$};
  \node[draw, carrier object, carrier fill, circle, minimum size=0.65cm, font=\scriptsize] (z5) at (2,2) {$z_5$};
  \node[draw, carrier object, carrier fill, minimum size=0.85cm] (g) at (3.4,1) {$g$};
  \node[draw, carrier object, circle, minimum size=0.9cm, font=\scriptsize] (zx) at (4.9,1) {$z(x)$};
  \foreach \zi in {z1,z2,z3,z4,z5} {
    \draw[carrier arrow, carrier muted] (\zi) -- (g.west);
  }
  \draw[carrier arrow] (g.east) -- (zx.west);
\end{tikzpicture}

  \caption{$A_2$: локальный закон $z(x,t{+}1)=g(\{z(y,t)\}_{y\in N(x)})$.}
  \label{fig:axiom-a2-local-law}
\end{figure}

\begin{axiom}[A3 · Унитарность / сохранение информации]
\label{ax:A3}
Норма $N(t)$ --- инвариант эволюции:
\[
N(t+1)=N(t)\quad\text{для всех }t.
\]
Линейная часть $g_{\mathrm{lin}}$ --- произведение локальных унитарных операторов на рёбрах $N(x)$.
Нелинейная часть --- только унитарная фаза $z\mapsto z\,e^{i\Phi(z)}$ на $\mathbb{C}^2$; модуль $|z|$ не диссипирует.
\end{axiom}

\begin{axiom}[A4 · U(1) заряд и SU(2) спин]
\label{ax:A4}
Состояние узла --- спинор $z=(z_1,z_2)^\top\in\mathbb{C}^2$.
Глобальная U(1)-фаза $z\to z\,e^{i\theta}$ --- симметрия вакуума.
Локальное столкновение реализуется SU(2)-вращателем $R(\Phi)$ на обеих компонентах.
\end{axiom}

\begin{figure}[htbp]
  \centering
  \begin{tikzpicture}[carrier panel]
  \draw[carrier muted, dashed] (0,0) circle (1);
  \draw[carrier arrow] (0,0) -- (35:1);
  \node at (35:1.15) {$z$};
  \draw[carrier arrow, carrier muted, dashed] (0,0) -- (95:1);
  \node[carrier muted] at (95:1.2) {$z\,e^{i\symCapitalPhi}$};
  \draw[carrier object] (35:0.38) arc (35:95:0.38);
  \node at (65:0.55) {$\symCapitalPhi$};
\end{tikzpicture}

  \caption{$A_3$: сохранение нормы $|z|$, меняется только фаза.}
  \label{fig:axiom-a3-phase}
\end{figure}

\begin{figure}[htbp]
  \centering
  \begin{tikzpicture}[carrier panel]
  \draw[carrier muted, dashed] (-1.1,0) circle (0.5);
  \draw[carrier muted, dashed] (1.1,0) circle (0.5);
  \draw[carrier arrow] (-1.1,0) -- (-35:0.42);
  \node at (-35:0.58) {$z_1$};
  \draw[carrier arrow] (-1.1,0) -- (112:0.33);
  \node at (112:0.48) {$z_2$};
  \draw[carrier arrow, carrier muted, dashed] (1.1,0) -- (-11:0.42);
  \node[carrier muted] at (-11:0.58) {$z'_1$};
  \draw[carrier arrow, carrier muted, dashed] (1.1,0) -- (136:0.33);
  \node[carrier muted] at (136:0.48) {$z'_2$};
  \draw[carrier arrow] (-0.55,0) -- (0.55,0);
  \node[anchor=south] at (0,0.12) {$R(\symCapitalPhi)$};
\end{tikzpicture}

  \caption{$A_4$: $SU(2)$-вращатель $R(\Phi)$ на $\mathbb{C}^2$.}
  \label{fig:axiom-a4-su2}
\end{figure}

